\documentclass[10pt,a4paper]{amsart}
\usepackage[margin=19mm]{geometry}
\usepackage{amsmath,amssymb,mathtools}
\usepackage[scr=rsfs,cal=euler]{mathalfa}
\usepackage{xfrac}
\usepackage[unicode,hidelinks,bookmarks=false]{hyperref}
\newcommand{\ie}{i.e.\ }
\newcommand{\cf}{cf.\ }
\newcommand{\resp}{respectively\ }
\newcommand{\Subsection}{\subsection}
\providecommand{\nobreakdash}{\nobreak\hbox{-}}
\setlength{\emergencystretch}{2em}
\multlinegap0pt
\usepackage{amssymb}
\usepackage{enumerate}
\usepackage{mathtools}
\usepackage{tikz}
\usepackage{xcolor}
\usepackage{float}
\newtheorem{theorem}{Theorem}
\title{New Type degenerate Simsek numbers and related aspects}
\author{Lahcen Oussi}
\date{Source: arXiv:2508.11720v2 (CC BY 4.0)}
\begin{document}
\maketitle

\noindent\emph{Source and licence.} New Type degenerate Simsek numbers and related aspects. Version arXiv:2508.11720v2 (CC BY 4.0), distributed by arXiv under the Creative Commons Attribution 4.0 International licence, http://creativecommons.org/licenses/by/4.0/, as recorded in the arXiv licence metadata for this exact version and shown on the abstract page https://arxiv.org/abs/2508.11720v2. Full licence notice: \texttt{licenses/arxiv-2508.11720-CC-BY-4.0.txt}.\par\smallskip

\noindent\textit{Editorial scope.} Counted unit: the article's theorem thm:formntdSn (source0.tex lines 230--239). The article's complete original proof of it is delivered with it (source0.tex lines 240--257, one \texttt{proof} environment of 18 lines). No mathematics is written, repaired or re-derived: the statement and the proof are the article's own lines, in the article's own notation, and no retained line is substituted or re-flowed.  Retained uncounted as the source-local context the proof needs: lines 103--105; lines 191--193.  Nothing was omitted from any retained span, so the package carries no omission record.  No retained line contains a citation, so the wrapper carries no bibliography and no bibliography entry.  No retained line contains a figure environment, an image-inclusion command or drawing code, so no image is delivered and the package declares no asset dependency.  Editorial formatting only, in addition: an AMS \texttt{amsart} preamble that restates the article's own declarations and macros used by the retained lines, namely the environments \texttt{theorem} on separate counters, together with the standard packages those lines need.  The retained environments print in this document as Theorem 1 in order of appearance, whereas the article prints the same items with the numbering of its own full document.  The complete native source of the article as distributed in the arXiv e-print for this exact version is bundled unchanged as \texttt{sources/183346/source0.tex}.
\begin{equation}\label{SNFK}
    (x)_n=\sum_{k=0}^{n}S_{1}(n,k)x^k,
\end{equation}

 \begin{equation}\label{newgSn}
F_{k}(t; \alpha, \lambda):=\frac{(\lambda e^t+1)_{k,\alpha}}{k!}=\sum_{n=0}^{\infty}y_{1,\alpha}^{*}(n,k;\lambda)\frac{t^n}{n!}.  
 \end{equation}

\begin{theorem}\label{thm:formntdSn}
The numbers $y_{1, \alpha}^{*}(n,k;\lambda)$ can be expressed explicitly  as follows:
 \begin{equation}\label{expforntsn}
 y_{1,\alpha}^{*}(n,k;\lambda)=\frac{1}{k!}\sum_{\ell=0}^{k}\sum_{j=0}^{\ell}\binom{\ell}{j}\alpha^{k-\ell}S_{1}(k,\ell)\lambda^{j}j^{n}
    \end{equation}
or
\begin{equation}\label{expforntsn2}
y_{1, \alpha}^{*}(n,k;\lambda)=\frac{1}{k!}\sum_{\ell=0}^{k}\sum_{j=0}^{\ell}\binom{k}{\ell}\alpha^{\ell-j}\lambda^{j}j^{n}(1)_{k-\ell, \alpha}S_{1}(\ell,j).
\end{equation}
\end{theorem}

\begin{proof}
Using \eqref{SNFK}, we have 
\begin{align*}
\frac{(\lambda e^t+1)_{k,\alpha}}{k!}&=\frac{\alpha^{k}}{k!}\Big(\frac{\lambda e^t+1}{\alpha}\Big)_{k}\\
&=\frac{1}{k!}\sum_{\ell=0}^{k}\alpha^{k-\ell}S_{1}(k,\ell)(\lambda e^t+1)^{\ell}\\
&=\frac{1}{k!}\sum_{\ell=0}^{k}\alpha^{k-\ell}S_{1}(k,\ell)\sum_{j=0}^{\ell}\binom{\ell}{j}\lambda^{j}e^{jt}\\
&=\sum_{n=0}^{\infty}\Bigg(\frac{1}{k!}\sum_{\ell=0}^{k}\sum_{j=0}^{\ell}\alpha^{k-\ell}S_{1}(k,\ell)\binom{\ell}{j}\lambda^{j}j^n\Bigg)\frac{t^n}{n!}.
\end{align*}
Hence, by \eqref{newgSn}, we obtain the desired result.

To prove Eq.~\eqref{expforntsn2}, we use the fact that $(x+y)_{k, \alpha}=\sum\limits_{j=0}^{k}\binom{k}{j}(x)_{j,\alpha}(y)_{k-j, \alpha}$. Then,
\begin{align*}
\frac{\big(\lambda e^t+1\big)_{k, \alpha}}{k!}&=\frac{1}{k!}\sum_{\ell=0}^{k}\binom{k}{\ell}(\lambda e^t)_{\ell, \alpha}(1)_{k-\ell, \alpha}\\
&=\frac{1}{k!}\sum_{\ell=0}^{k}\binom{k}{\ell}\alpha^{\ell}\Big(\frac{\lambda e^t}{\alpha}\Big)_{\ell}(1)_{k-\ell, \alpha}\\
&=\sum_{n=0}^{\infty}\left(\frac{1}{k!}\sum_{\ell=0}^{k}\sum_{j=0}^{\ell}\binom{k}{\ell}\alpha^{\ell-j}\lambda^{j}j^{n}(1)_{k-\ell, \alpha}S_{1}(\ell,j)\right)\frac{t^n}{n!}.
\end{align*}
Hence the result follows.
\end{proof}

\end{document}
